\documentclass{article}

\textwidth=125mm
\textheight=200mm
\linespread{1.04}


\title{\textsc{An equivalence between Hedge and Mirror Descent}}
\usepackage{marginnote}
\reversemarginpar % switch default margin to the left
\usepackage{framed}
\usepackage{xcolor}
\usepackage{hyperref} % add backlinks to the end of bibliography items
\usepackage{nicefrac}
\usepackage{pgfplots}
\usepackage{floatrow}
\usepackage{amsmath}
\usepackage{amsthm}
\usepackage{listings}
\usepackage{tikz}
\usepackage{amssymb}
\usepackage[plain]{algorithm}
\lstset{emph={try,catch,throw,def,or,to,for,then,end,for,if,args,else,return},emphstyle=\color{blue!100!black}\textbf}
\lstset{mathescape=true,escapechar=\&,numbers=left,xleftmargin=3mm,linewidth=0.47\textwidth,framextopmargin=0pt,framexbottommargin=0pt,aboveskip=-1pt,belowskip=-1pt}
\usepackage{mdframed}
\usepackage{bm}
\usepackage[capitalise]{cleveref}
\usepackage[bf]{caption}
\usepackage{graphicx}
\usepackage{natbib}
\usepackage{xspace}

\definecolor{dkblue}{cmyk}{1,.54,.04,.19} 
\theoremstyle{plain}
\newtheorem{theorem}{Theorem}
\newtheorem{optprob}[theorem]{Optimisation Problem}
\newtheorem{lemma}[theorem]{Lemma}
\newtheorem{proposition}[theorem]{Proposition}
\newtheorem{corollary}[theorem]{Corollary}
\newtheorem{claim}[theorem]{Claim}
\theoremstyle{definition}
\newtheorem{definition}[theorem]{Definition}
\newtheorem{example}[theorem]{Example}
\newtheorem{remark}[theorem]{Remark}
\newtheorem{question}[theorem]{Question}
\newtheorem{fact}[theorem]{Fact}
\newtheorem{assumption}[theorem]{Assumption}



\newcommand{\ceil}[1]{\left\lceil #1 \right\rceil}
\newcommand{\R}{\mathbb R}
\newcommand{\Rd}{\mathbb{R}^d}
\newcommand{\argmin}{\operatornamewithlimits{arg\,min}}
\newcommand{\argmax}{\operatornamewithlimits{arg\,max}}
\newcommand{\ip}[1]{\langle #1 \rangle}
\newcommand{\Reg}{\mathfrak{R}}
\newcommand{\BReg}{\mathfrak{BR}}
\newcommand{\norm}[1]{\Vert #1 \Vert}
\newcommand{\E}{\mathbb E}
\newcommand{\PP}{\mathbb P}
\newcommand{\Var}{\mathbb V}
\newcommand{\cE}{\mathcal E}
\newcommand{\cK}{\mathcal K}
\newcommand{\cH}{\mathcal H}
\newcommand{\cF}{\mathcal F}
\newcommand{\cO}{\mathcal O}
\newcommand{\cS}{\mathcal S}
\newcommand{\cC}{\mathcal C}
\newcommand{\cL}{\mathcal L}
\newcommand{\cG}{\mathcal G}
\newcommand{\cA}{\mathcal A}
\newcommand{\cD}{\mathcal D}
\newcommand{\cP}{\mathcal P}
\newcommand{\BS}{\mathbb S}
\newcommand{\zeros}{ \bm 0}
\newcommand{\bbP}{\mathbb P}
\newcommand{\polylog}{\operatorname{polylog}}
\newcommand{\const}{\operatorname{const}}
\newcommand{\interior}{\operatorname{int}}
\newcommand{\conv}{\operatorname{conv}}
\newcommand{\diam}{\operatorname{diam}}
\renewcommand{\d}[1]{\operatorname{d}\!#1}
\newcommand{\abovelabel}[1]{\stackrel{\raisebox{0.5mm}{\textrm{\tiny \color{red} #1}}}}
\newcommand{\ta}{\tilde{a}}
\newcommand{\tP}{P'}
\newcommand{\tA}{A'}
\newcommand{\fmd}{\texttt{MD}}
\newcommand{\tfmd}{\widetilde{\texttt{MD}}}

\newcommand{\mdalg}{\textsf{MD}\xspace}
\newcommand{\smdalg}{\textsf{SMD}\xspace}
\newcommand{\expbyopt}{\textsf{ExpByOpt}\xspace}
\newcommand{\olo}{\texttt{OLO}\xspace}
\newcommand{\blo}{\texttt{BLO}\xspace}
\newcommand{\bco}{\texttt{BCO}\xspace}
\newcommand{\oco}{\texttt{OCO}\xspace}
\newcommand{\hl}{\hat{\ell}\xspace}


\begin{document}
\maketitle

\section{Simplex is a vector space}
The probability simplex $\Delta_d = \{x \in \R^d: x_i \geq 0, \sum_{i=1}^d x_i = 1\}$ is a vector space with the following operations:
\begin{itemize}
    \item Addition: $x \oplus y = \frac{(x_1y_1, \ldots, x_dy_d)}{\sum_{i=1}^{d} x_iy_i}$.
    \item Scalar multiplication: For $\alpha \in \R$, $\alpha \odot x = \frac{(x_1^\alpha, \ldots, x_d^\alpha)}{\sum_{i=1}^{d} x_i^\alpha}$.
    \item Zero element is the uniform distribution, e.g. $\textbf{0} = (\frac{1}{d}, \ldots, \frac{1}{d})$.
\end{itemize}
For an arbitrary vector $v \in \R^d$, define
\begin{align*}
    \overline{v} & = \frac{(\exp(-v_1), \ldots, \exp(-v_d))}{\sum_{i=1}^{d} \exp(-v_i)} \in \Delta_d\,.
\end{align*}
Note that $\overline{v+w} = \overline{v} \oplus \overline{w}$ and $\overline{\alpha v} = \alpha \odot \overline{v}$. Also note that $\overline{0} = \textbf{0}$.
% \begin{remark}
%     The use of $\odot$ and $\oplus$ here is just for being explicit.
%     Technically, we could use $+, \times, \sigma$ instead.
%     One just needs to use their definition when applied to elements of $\Delta_d$.
% \end{remark}

\section{Hedge is a linear function}
The action of mirror descent with negative entropy on the simplex (Hedge),
with learning rate $\eta > 0$ and loss vector $\ell_t \in \R^d$ at time $t$, is given by
\begin{align*}
    \fmd(\ell_{1:t}; \eta) & = \frac{
        \left(
        \exp(-\eta \sum_{s=1}^t \ell_{s,1}), \ldots, \exp(-\eta \sum_{s=1}^t \ell_{s,d})
        \right)
    }{\sum_{i=1}^{d} \exp(-\eta \sum_{s=1}^t \ell_{s,i})}
\end{align*}
where $\ell_{s,i}$ is the $i$-th coordinate of $\ell_s$.
Using the vector space operations defined above, we can rewrite this as
\begin{align*}
    \fmd(\ell_{1:t}; \eta) & = -\eta \odot \left( \bigoplus_{s=1}^t \overline{\ell_s} \right),
\end{align*}
which shows that Hedge is a linear function of the projected loss vectors $\overline{\ell_1}, \ldots, \overline{\ell_t}$.

\section{The Exp3 algorithm}
In the bandit setting, the learner only observes the loss of the action played at time $t$, which we denote by
\begin{align*}
    Y_t & = \ip{\ell_t, A_t}\,,
\end{align*}
where $A_t \sim \pi_t \in \Delta_d$ is the action played at time $t$.
The Exp3-type algorithm uses the estimate
\begin{align*}
    \hl_{t} & = \frac{Y_t}{\pi_{t,i}} \mathbf{1}_{A_t}\,,
\end{align*}
where $\mathbf{1}_{A_t}$ is the one-hot vector with a $1$ in the coordinate corresponding to $A_t$.
It is easy to see that
\begin{align*}
    \E_t[\hl_t] & = \ell_t\,.
\end{align*}

\section{Exp3 is Hedge, in expectation}
First, note that for a \textbf{random} element $X \in \mathcal{X} \subset \R^d$, we have
\begin{align*}
    \overline{\E}[\overline{X}] & = \bigoplus_{x \in \mathcal{X}} \PP(X = x) \odot \overline{x}\,.
\end{align*}
Note that
\begin{align*}
    \overline{\E[X]} = \overline{\E}[\overline{X}]\,.
\end{align*}
Using this, we can show that Exp3 is Hedge in expectation:
\begin{align*}
    \overline{\E}[\fmd(\hl_{1:t}; \eta)] & =
    \overline{\E}\left[-\eta \odot \left( \bigoplus_{s=1}^t \overline{\hl_s} \right)\right]                               \\
                                         & = -\eta \odot \left( \bigoplus_{s=1}^t \overline{\E}[\overline{\hl_s}] \right) \\
                                         & = -\eta \odot \left( \bigoplus_{s=1}^t \overline{\E[\hl_s]} \right)            \\
                                         & = -\eta \odot \left( \bigoplus_{s=1}^t \overline{\ell_s} \right)               \\
                                         & = \fmd(\ell_{1:t}; \eta)\,.
\end{align*}
However, note that this is not $\E[\fmd(\hl_{1:t}; \eta)] = \fmd(\ell_{1:t}; \eta)$.

\clearpage
\subsection{Imitated \mdalg for \blo}
Lemma~\ref{lem:1-sample-estimator} shows that it is possible to construct
a sequence
\begin{align*}
    \hl_1, \hl_2, \cdots, \hl_n\,,
\end{align*}
such that $\E[\hl_t] = \ell_t$.
A natural question is whether it is possible to construct a sequence
\begin{align*}
    \hat{a}_1, \hat{a}_2, \cdots, \hat{a}_n\,,
\end{align*}
such that $\E[\hat{a}_t] = a_t$, where $a_t = \fmd(\ell_{1:t-1})$.
If this is possible, then the \textbf{expected} regret of \smdalg is exactly the same as the regret of \mdalg.
This boils down to the following requirement on $P_t$:
\begin{framed}
    Choose a policy $P_t \in \cP(\cA)$ given $Y_{1:t-1}$ and $A_{1:t-1}$ such that
    $\E[A_t] = \fmd(\ell_{1:t-1})$.
\end{framed}
Couple of remarks:
\begin{itemize}
    \item The requirement $\E[A_t] = \fmd(\ell_{1:t-1})$ is \textbf{less demanding} than $\E_t[A_t] = \fmd(\hl_{1:t-1})$.
    \item The support of $P_t$ needs to be in $\cA$.
    \item Policy $P_t$ needs to be exploratory enough to ensure that $Y_t$ contains enough information so that $P_{t+1}$ can be chosen to satisfy the requirement in the next round.
    \item Policy $P_t$ needs to be chosen in a way such that $\E[\E_{A_t \sim P_t}(A_t)] = a_t = \fmd(\ell_{1:t-1})$.
\end{itemize}

\subsubsection{If \mdalg is linear}
Consider the case that $\fmd$ is linear, i.e.
\begin{align*}
    \fmd(\alpha_1 \ell_{1:t-1} + \alpha_2 \ell_{1:t-1}') = \alpha_1 \fmd(\ell_{1:t-1}) + \alpha_2 \fmd(\ell_{1:t-1}')\,.
\end{align*}
This is never the case because of the non-linearity of the function $F$ and the projection into $\cA$, but is almost true when $F$ is quadratic.
Suppose that we choose $P_t$ such that $\E[\hl_t] = \ell_t$, as ensured by lemma~\ref{lem:1-sample-estimator},
and use $P_t = \delta_{\fmd(\hl_{1:t-1})}$, i.e. a Dirac delta at $\fmd(\hl_{1:t-1})$.
Then we have
\begin{align*}
    \E_t[A_t] & = \E_t[\fmd(\hl_{1:t-1})]                                      \\
              & = \fmd(\E_t[\hl_{1:t-1}]) \quad (\text{by linearity of } \fmd) \\
              & = \fmd(\ell_{1:t-1})\,,
\end{align*}
which is the stronger version of the requirement since $\E_t$-unbiasedness implies $\E$-unbiasedness.

\subsubsection{The possibility of playing $a_2$ in expectation}
Here, we are interested in the case of $\cA = \Delta(d)$ and $F$ being the negentropy.
One should ask if the requirement $\E[A_2] = \fmd(\ell_{1}; \eta)$ is possible to satisfy,
before worrying about the general case.

Note that $P_1$ is the uniform distribution over $[d]$.
So the task of the learner is to choose a distribution $P_2$ over $[d]$ such that
\begin{align*}
    \E[A_2] = \fmd(\ell_1; \eta) = \frac{\exp(-\eta \ell_{1,i})_{i=1}^{d}}{\sum_{j=1}^{d} \exp(-\eta \ell_{1,j})}\,.
\end{align*}
Note that $P_2$ is itself random, and is chosen based on $A_1$ and $Y_1 = \ell_{1, A_1}$.
Therefore, $P_2$ can be seen as the output of a function $D: [d] \times \R \to \Delta(d)$.
We have
\begin{align*}
    \E[A_2] = \E[\E_{A_2 \sim P_2}[A_2]] = \E[D(A_1, Y_1)] = \frac{1}{d} \sum_{i=1}^{d} D(i, \ell_{1,i})\,,
\end{align*}
and one can check that no such function $D$ exists.

\clearpage
\bibliographystyle{plainnat}
\bibliography{all}


\end{document}





